%% BRANCH project proposal, version 2 (wording revised 4 November 2026). IJCAI two column format.
%% Build: latexmk -pdf BRANCH_Project_Proposal_v2.tex
\documentclass{article}
\pdfpagewidth=8.5in
\pdfpageheight=11in
\usepackage{ijcai26}
\usepackage{times}
\usepackage{url}
\usepackage[hidelinks]{hyperref}
\usepackage[utf8]{inputenc}
\usepackage[small]{caption}
\usepackage{amsmath}
\usepackage{booktabs}
\usepackage{enumitem}
\usepackage[round,authoryear]{natbib}
\usepackage{titlesec}
\titlespacing*{\section}{0pt}{6pt}{3pt}
\titlespacing*{\subsection}{0pt}{4pt}{2pt}
\setlength{\bibsep}{0pt}
\setcitestyle{aysep={}}
\setlist{nosep,leftmargin=*}
\urlstyle{same}

\title{BRANCH: A BDI Agent without a Language Model that Completes an openLCA Life Cycle Assessment
from a Bill of Materials and a Location, Keeping the Human in Command}
\author{Ashira Benchimol
\affiliations
Department of Informatics Engineering, University of Coimbra\\
Project Proposal, version 2 (26 October 2026, wording revised on 4 November 2026),
Human Centered Artificial Intelligence, 2026/2027}

\begin{document}
\maketitle

\begin{abstract}
openLCA is the reference open source tool for life cycle assessment (LCA), but it calculates and does not
advise. Before the first result appears, the user must take every modelling choice that ISO 14044 and the
ILCD Handbook leave to expert judgement. I propose BRANCH, an agent for openLCA that contains no language
model. BRANCH is a BDI agent, which reasons with beliefs, desires and intentions. From a bill of materials and a location, it completes all remaining modelling
choices as explicit random variables and computes the result by Monte Carlo from impact factors that it
takes from openLCA. It asks the human only the questions whose answer can reverse the verdict, and it
chooses them with a decision flip query strategy derived from active learning. It returns an importable
openLCA model and a report structured by ISO 14044. A first prototype on real data reaches a certified
verdict with less than half the questions of uncertainty sampling. The project draws on four topics of
the course, namely autonomous agents, human and AI cooperation, human oversight and trustworthy AI.
\end{abstract}

\section{Problem and Motivation}
Life cycle assessment quantifies the environmental impacts of a product from raw material extraction to
the end of its life \citep{iso14040}. European regulation now asks companies for information of LCA
quality. The Packaging and Packaging Waste Regulation (EU) 2025/40 \citep{ppwr2024}, the Corporate
Sustainability Reporting Directive, the Environmental Footprint recommendation and the forthcoming Digital
Product Passport all require defensible and comparable environmental data for each product. Most of the
firms that must comply are small and medium enterprises without an LCA expert.

openLCA is free and widely used, reads every major inventory database and implements the Environmental
Footprint impact methods. Yet before the first number appears, the user must take dozens of
decisions that ISO 14044 and the ILCD Handbook leave to expert judgement \citep{ilcd2010,iso14040}. The
user chooses the functional unit, the decision context (Situation A, B or C), the system boundary and its
cut off, the allocation rule, the end of life model, the recycled content, the electricity mix of the
location and the transport scenario. The user also decides which of several near identical datasets
represents a material, for instance bottle grade PET in Portugal. Each decision is a branch in a tree, and only some
branches reverse the verdict between two alternatives.

The ILCD Handbook prescribes an iterative approach in which the analyst refines a rough model only where
it matters \citep[25]{ilcd2010}, but a non expert cannot know what matters. The outcome is either no LCA
at all or an LCA whose conclusion rests on defaults that the user never saw. The problem I address is therefore not computing an LCA, because openLCA already does that. The problem is
deciding, for a given product and place, which modelling choices a human must make and which ones data can
settle, and making that division explicit, auditable and cheap. This is a human centered AI problem,
because I look for the adequate level of automation while the human keeps control of the choices that
carry values.

\section{Background}
An LCA is only as good as the modelling choices behind it. ISO 14040 fixes the four phases and requires
that any comparison between products treat uncertainty and sensitivity explicitly \citep{iso14040}. The
ILCD Handbook makes those requirements operational and shows how many choices a practitioner must take
before computing anything \citep{ilcd2010}. Two papers in this course illustrate the stakes. Whether capital goods are inside the
boundary changes the result for some product groups and not for others \citep{frischknecht2007}, and the
same product can receive very different carbon footprints depending on boundary and allocation
\citep{schmidt2009}. The problem, then, is not calculating but choosing.

Current tools help only halfway. openLCA propagates uncertainty by Monte Carlo and ranks inputs by their
contribution to the variance of the result, but an input can be very uncertain and still never change
which product wins.

A BDI agent holds beliefs about the world, forms
desires, commits to intentions and reconsiders them as new information arrives
\citep{wooldridge2009,russell2020}. In BRANCH, the beliefs are the open modelling choices and the desire
is a robust verdict. At each step, the intention is to compute, to ask or to stop.

Deciding what to ask is an active learning problem. In active learning, the machine picks the instances
it most needs labelled, and it uses uncertainty sampling such as least confident, margin and entropy
\citep{settles2009}. The query strategy, together with the batch size, determines how little labelled
data suffices \citep{faria2022}. In BRANCH, the instance is a modelling choice and the label is the
human's decision. This design also keeps the human in command, because value laden choices stay with the
person and routine ones go to the agent. This division of control distinguishes human in the loop cooperation
from machine in the loop cooperation \citep{shneiderman2022,dafoe2020}.

Recent systems automate LCA with large language model agents that start from a product name or photo
\citep{zhang2026}; the modelling choices stay with the agent rather than with the user. To my knowledge,
no tool takes a bill of materials and a location, builds the openLCA model and asks only about the choices
that could change the conclusion.

\section{Goals}
The aim is to design, implement and evaluate BRANCH (Belief Revising Agent for Non expert Certified Human
in command LCA), a companion agent for openLCA with four abilities.
\begin{enumerate}
\item BRANCH accepts as input only a bill of materials (material type and mass of each component), the
location of production and use, and optionally a second alternative.
\item It completes every remaining ISO 14044 and ILCD modelling choice from data and codified rules, and
it keeps each choice as an explicit random variable instead of a hidden default.
\item It asks the human only the questions whose answer can flip the verdict, in order of decision value.
It accepts ``I don't know'' as an answer, and it stops when the verdict is robust.
\item It exports an importable openLCA product system and a report structured by ISO 14044.
\end{enumerate}
I will consider the project successful if three conditions hold.
\begin{enumerate}[label=(\roman*)]
\item On a benchmark of packaging cases, BRANCH reproduces openLCA results built by experts within an
agreed tolerance.
\item BRANCH reaches a certified verdict with fewer questions than the uncertainty sampling baselines.
\item At least 95 percent of certified verdicts hold under expert models that were built independently
of the agent's own priors.
\end{enumerate}
The project will contribute a decision flip query strategy for LCA, a BDI formalisation of the ILCD
iterative approach and an open source companion for openLCA, released under MPL 2.0 like openLCA.

\section{Approach}
\subsection{Beliefs}
BRANCH is a symbolic and probabilistic agent, and it contains no language model. It keeps a probability
distribution over every unstated modelling choice. Discrete choices, such as the allocation rule, the end
of life model, the decision context and the selection among candidate datasets, receive categorical
priors derived from ILCD and Environmental Footprint provisions. Continuous choices, such as recycled
content, collection and recycling rates, transport distance, reuse trips, washing energy and the
emission factor of electricity, receive priors centered on the statistics of the location. The spread of
these priors comes from the pedigree data of the database. The bill of materials and the location fix the
deterministic part of the belief base. BRANCH flags value laden choices, such as the inclusion of capital
goods, and always asks about them.

\subsection{Desires and intentions}
The desire is a verdict (the impact of A, or the sign of A minus B) whose probability of reversal stays
below 5 percent, obtained with the fewest questions.
A control loop in the style of Wooldridge chooses the intentions. The agent can compute, which means that
it propagates its beliefs by Monte Carlo. It can ask the human one question, or it can certify, which
means that it stops and reports. The agent reconsiders its plan whenever an answer changes the ranking of
the open choices.

\subsection{Query strategy}
For each open choice, the agent estimates by preposterior analysis how much the risk of a wrong verdict
falls once the human fixes that choice, and it selects the choice with the largest expected reduction.
This is active learning turned towards the human. The instance is a modelling choice and the label is the
human's decision. I will compare this strategy with least confident, margin and entropy sampling over the
same choices, with random order and with asking every question.

Version 1 measured that risk as $\min(p,1-p)$, where $p$ is the probability that A beats B. The prototype
showed that this measure makes the agent stall. When no single answer can carry $p$ across 0.5, every
question looks useless. I now use $p(1-p)$. Its expected reduction is the variance of $p$ across answers,
which is a sensitivity index of the verdict rather than of the impact.

\subsection{Outputs}
BRANCH writes the resolved model as openLCA JSON-LD, with its processes, flows, product systems and
parameters, so that the user can open, inspect and edit it in openLCA. JSON-LD is the exchange format of
openLCA itself and the cheapest way to hand a finished model back to it. The report lists the verdict, its estimated probability and the impacts of
both alternatives. It also lists every choice with its provenance (human, database, rule or ``don't
know'') and the smallest change in a single choice that reverses the verdict. When the comparison will be
disclosed to the public, the agent states at the start and in the report that ISO 14044 requires a
critical review panel. BRANCH prepares the study for that review but does not replace it.

\subsection{Existing implementations}
I use openLCA, olca-ipc and olca-schema as they are. LCA is linear in the amounts of each exchange, so
openLCA computes once, through olca-ipc, the impact of every candidate dataset per unit (per kg, kWh or
tonne kilometre). Monte Carlo then runs in Python on these factors, because the simulation built into
openLCA cannot be conditioned on partial answers. Each step takes milliseconds, so the person never waits.
The BDI loop and the query strategy are new code.

\section{Materials}
The software consists of openLCA 2.6, the olca-ipc and olca-schema Python packages, Python with NumPy
and SciPy, and GitHub. The inventory comes from the Life Cycle Inventory database of the Swiss Federal Administration,
BAFU:2026, which openLCA Nexus distributes for free \citep{bafu2026}. I combine it with the free openLCA
LCIA methods pack and its Environmental Footprint 3.1 method. I first planned to use the Environmental
Footprint database, but it is free only for registered Environmental Footprint studies and its licence
agreements expired in December 2025, so I replaced it. The terms of BAFU:2026 allow its use in calculators
but forbid redistribution, so each user extracts the factors on their own computer.

Eurostat provides the recycling rates of packaging waste,
ENTSO-E and the Brazilian National Energy Balance provide the emission factors of electricity, and the
Environmental Footprint rules provide the default transport distances. The benchmark uses the bills of
materials of beverage packaging (single use PET, refillable glass, aluminium cans and beverage cartons)
for Portugal, Germany and Brazil. I chose these products because published comparative LCAs exist for
them and their verdict is known to depend on reuse rates, recycled content and distance. The expert
models come from the ELCD PET bottle case study on openLCA Nexus and from the published studies.

\section{Evaluation}
\paragraph{Accuracy.} For each benchmark case, I will compute the relative error between BRANCH and the
expert model in every Environmental Footprint category, with bootstrap confidence intervals.

\paragraph{Question efficiency.} A simulated human answers from an expert model. Following
\citet{faria2022}, I will count the questions each strategy needs to reach a certified verdict over 30
runs and batch sizes from 1 to 5, and I will compare the strategies with paired Wilcoxon tests.

\paragraph{Calibration.} I will measure the share of certified verdicts that turn out wrong, with exact
binomial intervals, under two simulated humans. The first answers from the agent's own priors and the second
from their tails, so that the agent is not judged by its own assumptions.

\paragraph{Human study.} Three to five non experts and two practitioners will solve one case with openLCA
alone and one with BRANCH. I will report their time, their errors and their trust descriptively.

\subsection{Preliminary results}
The first test compares single use PET with refillable glass, per litre delivered in Portugal, for
climate change, with 14 open choices and factors from BAFU:2026 (Table~\ref{tab:prelim}). The decision
flip strategy needs less than half the questions of the course baselines (Wilcoxon $p<10^{-5}$), and it
always asks about the distribution distance first. Over 300 runs, certified verdicts were wrong in 1.0
percent of cases with correct priors (95\% interval 0.2 to 2.9) and in 4.3 percent with wrong priors
(2.3 to 7.3), which stays within the 5 percent target. Twelve of the fourteen priors are still
placeholders, so I do not claim an environmental verdict yet.

\begin{table}[h]
\vspace{-4pt}
\centering
\small
\begin{tabular}{lcc}
\toprule
Strategy & Correct priors & Wrong priors\\
\midrule
Decision flip, $p(1-p)$ & 4.6 & 3.2\\
Entropy, least confident & 10.5 & 9.2\\
Margin & 11.3 & 9.8\\
Random & 10.5 & 9.5\\
\bottomrule
\end{tabular}
\caption{Mean number of questions to a certified verdict over 30 runs.}
\label{tab:prelim}
\vspace{-8pt}
\end{table}

The paper will state three limitations. The agent covers attributional studies in ILCD Situation A,
its priors are only as good as the regional statistics behind them, and it does not replace the critical
review that ISO 14044 requires for public comparative assertions.

\begin{thebibliography}{99}
\scriptsize
\setlength{\itemsep}{0pt}
\bibitem[{BAFU(2026)}]{bafu2026}
BAFU (Federal Office for the Environment). 2026. \emph{Life Cycle Inventory Database of the Swiss Federal Administration, BAFU:2026}. Version 1.1. Bern: Federal Office for the Environment. Distributed via openLCA Nexus.
\bibitem[{Dafoe et~al.(2020)}]{dafoe2020}
Dafoe, Allan, Edward Hughes, Yoram Bachrach, Tantum Collins, Kevin R. McKee, Joel Z. Leibo, Kate Larson, and Thore Graepel. 2020. ``Open Problems in Cooperative AI.'' arXiv preprint arXiv:2012.08630. \url{https://arxiv.org/abs/2012.08630}.
\bibitem[{European Commission(2010)}]{ilcd2010}
European Commission, Joint Research Centre, Institute for Environment and Sustainability. 2010. \emph{International Reference Life Cycle Data System (ILCD) Handbook: General Guide for Life Cycle Assessment, Detailed Guidance}. EUR 24708 EN. Luxembourg: Publications Office of the European Union.
\bibitem[{European Parliament and Council(2024)}]{ppwr2024}
European Parliament and Council. 2024. ``Regulation (EU) 2025/40 of 19 December 2024 on Packaging and Packaging Waste, Amending Regulation (EU) 2019/1020 and Directive (EU) 2019/904, and Repealing Directive 94/62/EC.'' \emph{Official Journal of the European Union} L, 22 January 2025. \url{http://data.europa.eu/eli/reg/2025/40/oj}.
\bibitem[{Faria et~al.(2022)}]{faria2022}
Faria, Bruno, Dylan Perdig\~ao, Joana Br\'as, and Lu\'is Macedo. 2022. ``The Joint Role of Batch Size and Query Strategy in Active Learning-Based Prediction: A Case Study in the Heart Attack Domain.'' In \emph{Progress in Artificial Intelligence: EPIA 2022}. Cham: Springer.
\bibitem[{Frischknecht et~al.(2007)}]{frischknecht2007}
Frischknecht, Rolf, Hans-J\"org Althaus, Christian Bauer, Gabor Doka, Thomas Heck, Niels Jungbluth, Daniel Kellenberger, and Thomas Nemecek. 2007. ``The Environmental Relevance of Capital Goods in Life Cycle Assessments of Products and Services.'' \emph{International Journal of Life Cycle Assessment} 12 (Special Issue 1): 7--17.
\bibitem[{ISO(2006)}]{iso14040}
ISO. 2006. \emph{ISO 14040:2006 Environmental Management: Life Cycle Assessment, Principles and Framework}. Geneva: International Organization for Standardization.
\bibitem[{Russell and Norvig(2020)}]{russell2020}
Russell, Stuart, and Peter Norvig. 2020. \emph{Artificial Intelligence: A Modern Approach}. 4th ed. Hoboken: Pearson.
\bibitem[{Schmidt(2009)}]{schmidt2009}
Schmidt, Mario. 2009. ``Carbon Accounting and Carbon Footprint: More than Just Diced Results?'' \emph{International Journal of Climate Change Strategies and Management} 1 (1): 19--30.
\bibitem[{Settles(2009)}]{settles2009}
Settles, Burr. 2009. \emph{Active Learning Literature Survey}. Computer Sciences Technical Report 1648. Madison: University of Wisconsin--Madison.
\bibitem[{Shneiderman(2022)}]{shneiderman2022}
Shneiderman, Ben. 2022. \emph{Human-Centered AI}. Oxford: Oxford University Press.
\bibitem[{Wooldridge(2009)}]{wooldridge2009}
Wooldridge, Michael. 2009. \emph{An Introduction to MultiAgent Systems}. 2nd ed. Chichester: Wiley.
\bibitem[{Zhang et~al.(2026)}]{zhang2026}
Zhang, Zhihan, et~al. 2026. ``Sustainability Assessment Using Multimodal AI Agents.'' \emph{Nature Electronics}. \url{https://doi.org/10.1038/s41928-026-01653-w}.
\item[] \textit{Use of AI tools.} I used Claude (Anthropic) as a writing, research and coding assistant.
The design decisions, the problem and the final wording are my own, and I checked every reference
against its source.
\end{thebibliography}

\end{document}
